problem:  
Let \(G \subset \mathrm{SO}(n+1)\) act isometrically and with cohomogeneity one on the round sphere \(S^n\).  
Consider \(G\)-invariant hypersurfaces \(M_\gamma = G\cdot \gamma(t)\subset S^n\) determined by a profile curve \(\gamma: [0,L]\to S^n/G\).  
The area functional \(\mathcal{A}(M_\gamma)=\mathrm{Vol}(M_\gamma)\) and the Willmore functional  
\[
\mathcal{W}(M_\gamma)=\int_{M_\gamma} (|H_\gamma|^2+1)\, d\mu_\gamma
\]
both reduce, via the Hsiang–Lawson symmetry reduction, to one‑dimensional functionals \(\mathcal{A}_\mathrm{red}(\gamma)\) and \(\mathcal{W}_\mathrm{red}(\gamma)\) on the orbit‑space interval.

**Problem (Equivariant Willmore correspondence near minimal orbits):**  
For a fixed cohomogeneity‑one action \(G\curvearrowright S^n\), let \(M_0 \in \mathcal{M}_G\) be a \(G\)-invariant minimal hypersurface (a critical point of \(\mathcal{A}_\mathrm{red}\)).  
Let \(\mathcal{L}_\mathrm{min}\) and \(\mathcal{L}_\mathrm{will}\) denote the linearizations (Jacobi or stability operators) of \(\mathcal{A}_\mathrm{red}\) and \(\mathcal{W}_\mathrm{red}\) at the same profile \(\gamma_0\) defining \(M_0\).

Determine whether there exists any nontrivial cohomogeneity‑one action \(G\subset \mathrm{SO}(n+1)\) for which the following *spectral coincidence* holds:
\[
\mathrm{spec}(\mathcal{L}_\mathrm{min}) = \mathrm{spec}(\mathcal{L}_\mathrm{will})
\]
(in the sense of eigenvalue sets of the associated Sturm–Liouville or higher‑order operators on invariant variations), possibly up to a constant scaling factor.  
If no such action exists, characterize all actions for which the principal parts of these two linearized operators coincide.  

Why is it a "good" problem:  
This modified formulation bypasses the trivial and structurally impossible demand for exact coincidence of the minimal and Willmore ODEs, focusing instead on a *second‑variation correspondence* near symmetric minimal submanifolds.  The question probes whether some cohomogeneity‑one geometric symmetry can force the analytic stability behaviors of the two different variational functionals—area and bending energy—to align.  
Such a phenomenon, if it exists, would uncover a deep analytic harmony between first‑ and second‑order curvature variational problems that is not forbidden by general equation order considerations.  

Furthermore, determining when their linearized spectra coincide unites several areas:
- representation theory of \(G\) (dictating invariant variation spaces),
- spectral geometry of Hsiang–Lawson ODE operators,
- the analytic comparison of second‑ and fourth‑order geometric differential operators.  

Even proving non‑existence would reveal rigidity mechanisms tying group symmetry to curvature energy stability, while a positive example would suggest a new duality principle in symmetric submanifold theory.  
Its statement is compact and accessible—focused on the interplay of symmetry, geometry, and spectral analysis—yet it invites profound investigation, satisfying the “narrow entrance, deep depth” hallmark of a genuinely good research problem.